% Options for packages loaded elsewhere
% Options for packages loaded elsewhere
\PassOptionsToPackage{unicode}{hyperref}
\PassOptionsToPackage{hyphens}{url}
\PassOptionsToPackage{dvipsnames,svgnames,x11names}{xcolor}
%
\documentclass[
  11pt,
]{article}
\usepackage{xcolor}
\usepackage[margin=1in]{geometry}
\usepackage{amsmath,amssymb}
\setcounter{secnumdepth}{-\maxdimen} % remove section numbering
\usepackage{iftex}
\ifPDFTeX
  \usepackage[T1]{fontenc}
  \usepackage[utf8]{inputenc}
  \usepackage{textcomp} % provide euro and other symbols
\else % if luatex or xetex
  \usepackage{unicode-math} % this also loads fontspec
  \defaultfontfeatures{Scale=MatchLowercase}
  \defaultfontfeatures[\rmfamily]{Ligatures=TeX,Scale=1}
\fi
\usepackage{lmodern}
\ifPDFTeX\else
  % xetex/luatex font selection
\fi
% Use upquote if available, for straight quotes in verbatim environments
\IfFileExists{upquote.sty}{\usepackage{upquote}}{}
\IfFileExists{microtype.sty}{% use microtype if available
  \usepackage[]{microtype}
  \UseMicrotypeSet[protrusion]{basicmath} % disable protrusion for tt fonts
}{}
\makeatletter
\@ifundefined{KOMAClassName}{% if non-KOMA class
  \IfFileExists{parskip.sty}{%
    \usepackage{parskip}
  }{% else
    \setlength{\parindent}{0pt}
    \setlength{\parskip}{6pt plus 2pt minus 1pt}}
}{% if KOMA class
  \KOMAoptions{parskip=half}}
\makeatother
% Make \paragraph and \subparagraph free-standing
\makeatletter
\ifx\paragraph\undefined\else
  \let\oldparagraph\paragraph
  \renewcommand{\paragraph}{
    \@ifstar
      \xxxParagraphStar
      \xxxParagraphNoStar
  }
  \newcommand{\xxxParagraphStar}[1]{\oldparagraph*{#1}\mbox{}}
  \newcommand{\xxxParagraphNoStar}[1]{\oldparagraph{#1}\mbox{}}
\fi
\ifx\subparagraph\undefined\else
  \let\oldsubparagraph\subparagraph
  \renewcommand{\subparagraph}{
    \@ifstar
      \xxxSubParagraphStar
      \xxxSubParagraphNoStar
  }
  \newcommand{\xxxSubParagraphStar}[1]{\oldsubparagraph*{#1}\mbox{}}
  \newcommand{\xxxSubParagraphNoStar}[1]{\oldsubparagraph{#1}\mbox{}}
\fi
\makeatother


\usepackage{longtable,booktabs,array}
\newcounter{none} % for unnumbered tables
\usepackage{calc} % for calculating minipage widths
% Correct order of tables after \paragraph or \subparagraph
\usepackage{etoolbox}
\makeatletter
\patchcmd\longtable{\par}{\if@noskipsec\mbox{}\fi\par}{}{}
\makeatother
% Allow footnotes in longtable head/foot
\IfFileExists{footnotehyper.sty}{\usepackage{footnotehyper}}{\usepackage{footnote}}
\makesavenoteenv{longtable}
\usepackage{graphicx}
\makeatletter
\newsavebox\pandoc@box
\newcommand*\pandocbounded[1]{% scales image to fit in text height/width
  \sbox\pandoc@box{#1}%
  \Gscale@div\@tempa{\textheight}{\dimexpr\ht\pandoc@box+\dp\pandoc@box\relax}%
  \Gscale@div\@tempb{\linewidth}{\wd\pandoc@box}%
  \ifdim\@tempb\p@<\@tempa\p@\let\@tempa\@tempb\fi% select the smaller of both
  \ifdim\@tempa\p@<\p@\scalebox{\@tempa}{\usebox\pandoc@box}%
  \else\usebox{\pandoc@box}%
  \fi%
}
% Set default figure placement to htbp
\def\fps@figure{htbp}
\makeatother





\setlength{\emergencystretch}{3em} % prevent overfull lines

\providecommand{\tightlist}{%
  \setlength{\itemsep}{0pt}\setlength{\parskip}{0pt}}



 


\usepackage{amsmath,amssymb,enumitem,booktabs,tikz}
\usetikzlibrary{arrows.meta,positioning}
% tighter title block (same titling trick as the reference sheets) so Problem 1 fits page 1
\usepackage{titling}
\setlength{\droptitle}{-1.0in}
\posttitle{\par\end{center}}
\preauthor{}\postauthor{}\predate{}\postdate{}
\setlist{itemsep=2pt,topsep=3pt}
\usepackage{fancyhdr}\pagestyle{fancy}\fancyhf{}
\fancyhead[L]{\small DS701 Fall 2026}\fancyhead[R]{\small Homework 9}
\fancyfoot[C]{\small\thepage}
\renewcommand{\headrulewidth}{0.4pt}
\usepackage[breakable]{tcolorbox}
% --- answer-space macros (problems PDF only; no-ops in the solutions PDF) ---
% \answerspace{2in}  vertical room for handwritten work, sized from the solution
% \blank             a short fill-in rule; \blank[1.5in] for a wider one
% \answerpage        start the next problem on a fresh page (Gradescope page matching)
\newcommand{\answerspace}[1]{\par\vspace*{#1}}  % \vspace* so room is not dropped at a page break
\newcommand{\blank}[1][1.2in]{\rule{#1}{0.4pt}}
\newcommand{\answerpage}{\clearpage}
\newtcolorbox{solutionbox}{breakable,colback=blue!4,colframe=blue!45!black,title=Solution,fonttitle=\bfseries,left=4pt,right=4pt,top=3pt,bottom=3pt}
\makeatletter
\@ifpackageloaded{caption}{}{\usepackage{caption}}
\AtBeginDocument{%
\ifdefined\contentsname
  \renewcommand*\contentsname{Table of contents}
\else
  \newcommand\contentsname{Table of contents}
\fi
\ifdefined\listfigurename
  \renewcommand*\listfigurename{List of Figures}
\else
  \newcommand\listfigurename{List of Figures}
\fi
\ifdefined\listtablename
  \renewcommand*\listtablename{List of Tables}
\else
  \newcommand\listtablename{List of Tables}
\fi
\ifdefined\figurename
  \renewcommand*\figurename{Figure}
\else
  \newcommand\figurename{Figure}
\fi
\ifdefined\tablename
  \renewcommand*\tablename{Table}
\else
  \newcommand\tablename{Table}
\fi
}
\@ifpackageloaded{float}{}{\usepackage{float}}
\floatstyle{ruled}
\@ifundefined{c@chapter}{\newfloat{codelisting}{h}{lop}}{\newfloat{codelisting}{h}{lop}[chapter]}
\floatname{codelisting}{Listing}
\newcommand*\listoflistings{\listof{codelisting}{List of Listings}}
\makeatother
\makeatletter
\makeatother
\makeatletter
\@ifpackageloaded{caption}{}{\usepackage{caption}}
\@ifpackageloaded{subcaption}{}{\usepackage{subcaption}}
\makeatother
\usepackage{bookmark}
\IfFileExists{xurl.sty}{\usepackage{xurl}}{} % add URL line breaks if available
\urlstyle{same}
% fallback for those not using the hyperref driver hyperxmp:
\makeatletter
\@ifundefined{xmpquote}{\newcommand{\xmpquote}[1]{#1}}{}
\makeatother
\hypersetup{
  pdftitle={DS701 Homework 9: Causal Inference and Neural Networks},
  colorlinks=true,
  linkcolor={blue},
  filecolor={Maroon},
  citecolor={Blue},
  urlcolor={Blue},
  pdfcreator={LaTeX via pandoc}}


\title{DS701 Homework 9: Causal Inference and Neural Networks}
\usepackage{etoolbox}
\makeatletter
\providecommand{\subtitle}[1]{% add subtitle to \maketitle
  \apptocmd{\@title}{\par {\large #1 \par}}{}{}
}
\makeatother
\subtitle{Released Mon Nov 9 · due Mon Nov 16, 11:59 pm on Gradescope}
\author{}
\date{}
\begin{document}
\maketitle


\noindent\textbf{Name:}\ \rule{2.0in}{0.4pt}\qquad\textbf{BU ID:}\ \rule{1.2in}{0.4pt}
\vspace{6pt}

\textbf{How this works.} Answer \textbf{by hand on paper} (or a tablet
with a stylus), show your work, then scan and upload to Gradescope,
matching each problem to its pages. This is deliberate practice for the
closed-notes quizzes, which are in the same style --- so do it as much
as possible without a computer or an AI assistant. If you are stuck, you
can try AI in tutorial style or reach out to course staff. It is
\textbf{participation-graded}: a serious attempt at every problem earns
full credit, and worked solutions are released the day after the
deadline. Seven short parts; budget 1--2 hours.

\textbf{Covers:} L14 Causal Inference (Mon Nov 9), L16 Neural Networks
(Wed Nov 11). \textbf{Reference sheet:} quiz-3 sheet §14 and §16 ---
allowed and encouraged; every part is answerable from it plus what is
printed.

\subsection{Problem 1: Does the coupon work? (read the
table)}\label{problem-1-does-the-coupon-work-read-the-table}

\emph{Practices: crude vs.~stratified comparisons, Simpson's paradox,
and one standardization (L14).}

A streaming service e-mailed a ``welcome back'' coupon to some lapsed
subscribers: \(T=1\) if a user got the coupon, \(Y=1\) if they
resubscribed within 30 days, and \(Z=1\) if the account had been lapsed
\textbf{more than a year} (``long-lapsed''). Counts for 400 users; work
with resubscription \emph{rates} and the \textbf{risk difference}
\(P(Y{=}1\mid T{=}1)-P(Y{=}1\mid T{=}0)\) as the measure of effect:

\begin{center}\small\setlength{\tabcolsep}{7pt}
\begin{tabular}{@{}lcccc@{}}\toprule
 & \multicolumn{2}{c}{long-lapsed ($Z{=}1$)} & \multicolumn{2}{c}{recently lapsed ($Z{=}0$)}\\
\cmidrule(lr){2-3}\cmidrule(lr){4-5}
 & users & resubscribed & users & resubscribed\\ \midrule
coupon ($T{=}1$)    & 260 & 52 & 40 & 32\\
no coupon ($T{=}0$) &  40 &  4 & 60 & 42\\
\bottomrule\end{tabular}\end{center}

\begin{enumerate}
\def\labelenumi{(\alph{enumi})}
\tightlist
\item
  Compute the \textbf{crude} (pooled) resubscription rate for coupon and
  no-coupon users and the crude risk difference, then the risk
  difference \textbf{within} each stratum of \(Z\). Name the phenomenon,
  name the role \(Z\) plays, and point to the \textbf{two ratios in the
  table} that are its fingerprints.
\end{enumerate}

Crude rate, coupon: \blank[0.7in]\quad no coupon:
\blank[0.7in]\quad crude difference: \blank[0.7in]\\
Difference within \(Z{=}1\): \blank[0.7in]\quad within \(Z{=}0\):
\blank[0.7in]\answerspace{0.7in}

\begin{enumerate}
\def\labelenumi{(\alph{enumi})}
\setcounter{enumi}{1}
\tightlist
\item
  Use the sheet's adjustment (standardization) formula
  \(P(Y\mid do(T{=}t)) = \sum_z P(Y\mid T{=}t, Z{=}z)\,P(Z{=}z)\) to
  compute \(P(Y{=}1\mid do(T{=}1))\), \(P(Y{=}1\mid do(T{=}0))\) and
  their difference. In one sentence: what does this number estimate that
  the crude difference does not?
\end{enumerate}

\(P(Y{=}1\mid do(T{=}1))\) = \blank[0.6in]; \(P(Y{=}1\mid do(T{=}0))\) =
\blank[0.6in]; difference = \blank[0.6in]\answerspace{0.6in}

\answerpage

\subsection{Problem 2: One small DAG}\label{problem-2-one-small-dag}

\emph{Practices: reading a causal DAG --- fork vs.~collider --- choosing
the one variable to adjust for, and prediction vs.~intervention (L14).}

The same service wants the \textbf{total effect} of its
\textbf{autoplay} setting (\(T=1\) if on) on \textbf{weekly hours
watched} \(Y\). Its analysts agree on this DAG: long-tenured subscribers
(\(Z\) = tenure) are more likely to leave autoplay on \emph{and} watch
more regardless, and marketing keeps a list \(C\) of ``engaged users''
--- anyone with autoplay on \textbf{or} many hours watched. Because that
list has the cleanest data, one analyst proposes comparing autoplay-on
and autoplay-off users \textbf{within the list only}.

\begin{center}
\begin{tikzpicture}[>={Stealth[length=2.5mm]},
    var/.style={draw, circle, minimum size=7mm, inner sep=1pt, font=\small},
    lab/.style={font=\footnotesize, inner sep=1pt}]
  \node[var] (Z) at (2,0.8)  {$Z$};
  \node[var] (T) at (0,0)    {$T$};
  \node[var] (Y) at (4,0)    {$Y$};
  \node[var] (C) at (2,-0.8) {$C$};
  \draw[->] (Z) -- (T);  \draw[->] (Z) -- (Y);
  \draw[->] (T) -- (Y);
  \draw[->] (T) -- (C);  \draw[->] (Y) -- (C);
  \node[lab, right=3pt of Z] {tenure};
  \node[lab, left=3pt of T] {autoplay on};
  \node[lab, right=3pt of Y] {hours watched};
  \node[lab, right=3pt of C] {on the ``engaged users'' list};
\end{tikzpicture}
\end{center}

\begin{enumerate}
\def\labelenumi{(\alph{enumi})}
\tightlist
\item
  To estimate the total effect of \(T\) on \(Y\) there is exactly
  \textbf{one} variable in this DAG you should adjust for. Which one,
  and why --- name the building block (chain / fork / collider) that
  \(Z\) and \(C\) each form with \(T\) and \(Y\), and the sheet's rule
  for each. Then, in one or two sentences: what goes wrong with the
  ``within the list only'' proposal?
\end{enumerate}

Adjust for: \blank[0.6in]\qquad \(Z\) is a: \blank[1in]\qquad \(C\) is
a: \blank[1in]

\mbox{}

\par\nobreak\answerspace{1.1in}

\begin{enumerate}
\def\labelenumi{(\alph{enumi})}
\setcounter{enumi}{1}
\tightlist
\item
  \textbf{True or false, with one reason each} (a sentence is enough):
\end{enumerate}

\begin{enumerate}[label=(\roman*),leftmargin=3.2em,topsep=2pt]
\item A model that predicts hours watched very accurately from the autoplay
  setting and other variables therefore tells us how much hours watched would
  change if the service switched autoplay on for a user.
\item If the service randomly assigned autoplay on/off to new accounts, the
  plain difference in mean hours between the two arms would be an unbiased
  estimate of the average treatment effect, \emph{even though tenure $Z$ is not
  measured}.
\end{enumerate}

\noindent (i)\ \ True / False, because: \answerspace{0.65in}
\noindent (ii)\ True / False, because: \answerspace{0.65in}

\answerpage

\subsection{Problem 3: Training a small
network}\label{problem-3-training-a-small-network}

\emph{Practices: what the parameters of an MLP are, reading training and
validation loss, and one gradient-descent step (L16).}

A team fits a regression MLP with \textbf{6 (scaled) input features},
\textbf{one hidden layer of 8 ReLU units}, and \textbf{one linear
output} \(\hat y\), trained with mean squared error.

\begin{enumerate}
\def\labelenumi{(\alph{enumi})}
\tightlist
\item
  How many trainable parameters does the network have? Show the count as
  ``weights \(+\) biases'' for each layer (every entry of the sheet's
  \(\boldsymbol\Omega_k\) and \(\boldsymbol\beta_k\) counts). A teammate
  then suggests removing the ReLU (``the hidden layer is still there, so
  we keep the same capacity''). \textbf{True or false}, with one reason
  --- and say how many \emph{effective} parameters the ReLU-free network
  has.
\end{enumerate}

Total parameters: \blank[0.7in]\qquad ReLU-free claim is: True / False\\
Effective parameters without ReLU: \blank[0.7in]

\mbox{}

\par\nobreak\answerspace{1.1in}

\begin{enumerate}
\def\labelenumi{(\alph{enumi})}
\setcounter{enumi}{1}
\tightlist
\item
  Three training runs on the same data. Run A and Run B use the
  architecture above and differ only in \texttt{learning\_rate\_init}.
  Run C uses Run A's settings but the team forgot the
  \texttt{StandardScaler}, so feature 3 (annual income, in dollars)
  enters raw. Losses per epoch (Run C's row is in units of \(10^8\)):
\end{enumerate}

\begin{center}\small\setlength{\tabcolsep}{4.5pt}
\begin{tabular}{@{}lrrrrrrrrrr@{}}\toprule
epoch & 1 & 2 & 3 & 4 & 5 & 6 & 7 & 8 & 9 & 10\\ \midrule
Run A ($\eta=0.01$), train      & 0.92 & 0.61 & 0.44 & 0.35 & 0.29 & 0.25 & 0.22 & 0.19 & 0.17 & 0.15\\
Run A, validation               & 0.95 & 0.66 & 0.51 & 0.44 & 0.40 & \textbf{0.38} & 0.39 & 0.41 & 0.44 & 0.48\\
Run B ($\eta=1.0$), train       & 0.92 & 1.61 & 0.88 & 2.73 & 5.40 & 4.12 & 9.85 & 21.7 & 18.3 & 47.6\\
Run C (unscaled), train $/10^8$ & 2.1 & 2.0 & 2.0 & 1.9 & 1.9 & 1.9 & 1.9 & 1.9 & 1.9 & 1.9\\
\bottomrule\end{tabular}\end{center}

\begin{enumerate}[label=(\roman*),leftmargin=3.2em,topsep=2pt]
\item Diagnose \textbf{Run B} in one sentence, and name the knob and the
  direction you would turn it.
\item In \textbf{Run A}, what is happening from epoch 7 on, and what is it
  called? The team re-runs A with \texttt{early\_stopping=True} and
  \texttt{n\_iter\_no\_change=3}, which here means: \emph{stop as soon as the
  validation loss has failed to improve on its best value for 3 consecutive
  epochs, and keep the weights from the best epoch.} After which epoch does
  training stop, and which epoch's weights are kept?
\item Diagnose \textbf{Run C} in one sentence: why does one raw feature make
  gradient descent behave like this, and what is the fix?
\end{enumerate}

\noindent (i)\answerspace{1in}
\noindent (ii)\ \ this is called: \blank[1.2in]\quad stops after epoch: \blank[0.5in]\quad weights kept from epoch: \blank[0.5in]\answerspace{1in}
\noindent (iii)\answerspace{1in}

\begin{enumerate}
\def\labelenumi{(\alph{enumi})}
\setcounter{enumi}{2}
\tightlist
\item
  At some step of Run A, one hidden unit has incoming weights
  \(\mathbf w = (0.40,\,-0.90,\,1.50)\) and bias \(b = 0.20\), and
  backprop has already delivered the gradients
  \(\partial L/\partial \mathbf w = (-0.50,\,0.20,\,2.00)\) and
  \(\partial L/\partial b = -1.00\). Apply \textbf{one} gradient-descent
  update with \(\eta = 0.1\) and write the new \(\mathbf w\) and \(b\).
  Then: with \(\eta = 1.0\) instead, what would the third weight become,
  and what does the sheet say a learning rate that large risks?
\end{enumerate}

New \(\mathbf w\) = ( \blank[0.6in] , \blank[0.6in] ,
\blank[0.6in] )\qquad new \(b\) = \blank[0.6in]\qquad \(w_3\) with
\(\eta=1.0\): \blank[0.6in]

\mbox{}

\par\nobreak\answerspace{1.4in}




\end{document}
